% from ega2-1-fr.tex
\oldpage[II]{18}
($T$ indéterminée)~; c'est évident lorsque
$S=\operatorname{Spec}(\mathbf{Z})$, en vertu de
(\hyperref[II.1.7.6-fr]{1.7.6}), et on passe de là au cas général en
considérant le morphisme structural
$S\to\operatorname{Spec}(\mathbf{Z})$ et utilisant
(\hyperref[II.1.7.11-fr]{1.7.11, (iv)}). En raison de ce résultat, on pose
encore $\mathbf{V}(\mathcal{O}_S)=S[T]$, et on a donc la formule
\[
  S[T]=S\otimes_{\mathbf{Z}}\mathbf{Z}[T].
  \tag{1.7.12.1}\label{II.1.7.12.1-fr}
\]

L'identification du faisceau des germes de $S$-sections de $S[T]$ avec
$\mathcal{O}_S$, déjà vue dans
(\hyperref[I.3.3.15-fr]{I, 3.3.15}), se retrouve ici dans un contexte plus
général, comme cas particulier de (\hyperref[II.1.7.9-fr]{1.7.9}).
\end{env}

\begin{env}[1.7.13]
\label{II.1.7.13-fr}
Pour tout $S$-préschéma $X$, on a vu
(\hyperref[II.1.7.8-fr]{1.7.8}) que
$\operatorname{Hom}_S(X,S[T])$ s'identifie canoniquement à
$\operatorname{Hom}_{\mathcal{O}_S}(\mathcal{O}_S,\mathcal{A}(X))$, lequel
est canoniquement isomorphe à $\Gamma(S,\mathcal{A}(X))$, et est par suite
muni d'une structure d'anneau~; en outre, à tout $S$-morphisme $h:X\to Y$
correspond un morphisme
$\Gamma(\mathcal{A}(h)):\Gamma(S,\mathcal{A}(Y))
\to\Gamma(S,\mathcal{A}(X))$ pour les structures d'anneau
(\hyperref[II.1.1.2-fr]{1.1.2}). Quand on munit
$\operatorname{Hom}_S(X,S[T])$ de la structure d'anneau ainsi définie, on
voit donc qu'on peut considérer $\operatorname{Hom}_S(X,S[T])$ comme un
\emph{foncteur contravariant en $X$}, de la catégorie des $S$-préschémas
dans celle des anneaux. D'autre part,
$\operatorname{Hom}_S(X,\mathbf{V}(\mathcal{E}))$ s'identifie de même à
$\operatorname{Hom}_{\mathcal{O}_S}(\mathcal{E},\mathcal{A}(X))$ (où
$\mathcal{A}(X)$ est considéré comme un
\emph{$\mathcal{O}_S$-Module})~; on peut par suite le munir canoniquement
d'une structure de \emph{module} sur l'anneau
$\operatorname{Hom}_S(X,S[T])$, et on voit comme ci-dessus que le couple
\[
  \bigl(\operatorname{Hom}_S(X,S[T]),
  \operatorname{Hom}_S(X,\mathbf{V}(\mathcal{E}))\bigr)
\]
est un foncteur contravariant en $X$, à valeurs dans la catégorie dont les
éléments sont les couples $(A,M)$ formés d'un anneau $A$ et d'un $A$-module
$M$, les morphismes étant les di-homomorphismes.

Nous interpréterons ces faits en disant que $S[T]$ est un
\emph{$S$-schéma d'anneaux} et que $\mathbf{V}(\mathcal{E})$ est un
\emph{$S$-schéma de modules} sur le $S$-schéma d'anneaux $S[T]$
(cf. chap.~0, \textsection8).
\end{env}

\begin{env}[1.7.14]
\label{II.1.7.14-fr}
Nous allons voir que la structure de $S$-schéma de modules définie sur le
$S$-schéma $\mathbf{V}(\mathcal{E})$ permet de reconstituer le
$\mathcal{O}_S$-Module $\mathcal{E}$ à un isomorphisme unique près~: pour
cela, nous allons montrer que $\mathcal{E}$ est canoniquement isomorphe à un
sous-$\mathcal{O}_S$-Module de
$\mathbf{S}(\mathcal{E})=\mathcal{A}(\mathbf{V}(\mathcal{E}))$, défini au
moyen de cette structure. En effet
(\hyperref[II.1.7.4-fr]{1.7.4}) l'ensemble
$\operatorname{Hom}_{\mathcal{O}_S}(\mathbf{S}(\mathcal{E}),
\mathcal{A}(X))$ des homomorphismes de
\emph{$\mathcal{O}_S$-Algèbres} s'identifie canoniquement à
$\operatorname{Hom}_{\mathcal{O}_S}(\mathcal{E},\mathcal{A}(X))$, ensemble
d'homomorphismes de \emph{$\mathcal{O}_S$-Modules}~: si $h$, $h'$ sont deux
éléments de ce dernier ensemble, $s_i$ ($1\leq i\leq n$) des sections de
$\mathcal{E}$ au-dessus d'un ouvert $U\subset S$, $t$ une section de
$\mathcal{A}(X)$ au-dessus de $U$, on a par définition
\[
  (h+h')(s_1s_2\cdots s_n)
  =\prod_{i=1}^n\bigl(h(s_i)+h'(s_i)\bigr)
\]
et
\[
  (t.h)(s_1s_2\cdots s_n)=t^n\prod_{i=1}^n h(s_i).
\]

Cela étant, si $z$ est une section de $\mathbf{S}(\mathcal{E})$ au-dessus
de $U$, $h\mapsto h(z)$ est une application de
$\operatorname{Hom}_S(X,\mathbf{V}(\mathcal{E}))
=\operatorname{Hom}_{\mathcal{O}_S}(\mathbf{S}(\mathcal{E}),
\mathcal{A}(X))$ dans $\Gamma(U,\mathcal{A}(X))$. Nous allons
